\section{A3C: A Second Algorithm, a Decomposition and a Generative Market Model}
\label{sec:update-a3c}

\paragraph{Where the question comes from.} The original notebook and the re-run
(Section~\ref{sec:update-a2c}) use \emph{A2C}, i.e.\ synchronous learning. The obvious
cross-check is \emph{A3C}: several workers collect rollouts in parallel and push their
gradients into a shared network (Hogwild style, only the optimiser step is locked).
Environment, windows, episode seeds and evaluation functions are left unchanged, so the
comparison measures the learning algorithm and not the surrounding machinery. The final
evaluation again runs over \textbf{20 seeds (100--119)}. One peculiarity up front: A3C was
built only \emph{after} the A2C part was finished, so there is no notebook template against
which the numbers could be mirrored.

\subsection{The Baseline Has No Edge}

\begin{table}[H]
  \centering
  \caption{A3C on Boeing, 20 seeds. Margin $=$ agent $-$ buy-and-hold in NAV points
           (start $=1.0$); $p$ one-sided from the $t$-test ``margin $>0$''.
           ``R/R'' is the return/risk ratio in the test window.}
  \label{tab:update-a3c-arms}
  \small
  \begin{tabular}{lrrrrrr}
    \toprule
    Arm & Training & Test & 95\,\% CI test & Hit rate test & $p$ & R/R test\\
    \midrule
    Baseline                 & $+8.43$ & $+0.019$ & $[-0.06;+0.10]$ & 51.0\,\% & $0.32$  & $-0.24$\\
    \texttt{ent\_coef} 0.02  & $+8.37$ & $+0.047$ & $[-0.03;+0.12]$ & 64.0\,\% & $0.10$  & $-0.37$\\
    \texttt{obs\_noise} 0.01 & $+8.79$ & $+0.038$ & $[-0.03;+0.11]$ & 57.2\,\% & $0.14$  & $-0.45$\\
    \texttt{nstep} 10        & $+6.71$ & $+0.125$ & $[+0.01;+0.24]$ & 67.8\,\% & $0.019$ & $-0.05$\\
    Bundle (all three)       & $+7.69$ & $+0.145$ & $[+0.03;+0.26]$ & 70.2\,\% & $0.007$ & $-0.03$\\
    Bundle $+$ bootstrap     & $+0.02$ & $+0.218$ & $[+0.09;+0.35]$ & 77.8\,\% & $0.0012$& $+0.50$\\
    \texttt{nstep} 10 $+$ bootstrap & $+0.02$ & $+0.218$ & $[+0.09;+0.35]$ & 72.5\,\% & $0.0010$& $+0.38$\\
    \bottomrule
  \end{tabular}
\end{table}

The baseline repeats the A2C result: in the test window the margin is $+0.019$ NAV points,
the agent wins in \textbf{51.0\,\%} of the episodes --- a coin flip --- and the one-sided
$p$-value is $0.32$. \textbf{An unmodified A3C agent has no demonstrable edge over
buy-and-hold.} In the training window, at the same time, the margin is $+8.43$: the agent
has memorised the one historical path.

\subsection{Decomposition: Which Factor Carries the Effect?}

The obvious move is a bundle of three changes: \texttt{nstep} $5\to10$ (n-step returns),
\texttt{ent\_coef} $0.01\to0.02$ (entropy bonus) and \texttt{obs\_noise} $0\to0.01$
(Gaussian noise on the observation during training). The bundle lifts the test margin to
$+0.145$ ($p=0.007$). The only question is: \emph{which} of the three factors does that?

Each factor was therefore run on its own, again with 20 seeds (rows 2--4 of
Table~\ref{tab:update-a3c-arms}). The result is unambiguous and was surprising to me:

\paragraph{\texttt{nstep} is the active ingredient.} On its own, \texttt{nstep} $10$
already reaches $+0.125$ ($p=0.019$, hit rate 67.8\,\%) --- that is \textbf{85\,\%} of the
bundle's effect. The paired test bundle against \texttt{nstep} alone gives $+0.020$ at
$p=0.81$: statistically \textbf{indistinguishable}. The bundle is therefore not a sum of
levers but essentially \texttt{nstep} plus noise.

\paragraph{\texttt{ent\_coef} and \texttt{obs\_noise} are inert.} On their own they yield
$+0.047$ ($p=0.64$) and $+0.038$ ($p=0.72$) against the baseline, with 10 and 9 of 20
positive seed differences respectively --- a coin flip. They do not act on the overfitting
either: the in-sample margin stays at $+8.37$ and $+8.79$ against $+8.43$ for the baseline
($p=0.90$ and $0.53$). Only \texttt{nstep} $10$ lowers it significantly to $+6.71$
($p=0.0015$) --- and at the same time raises the test value.

\paragraph{Additivity.} The sum of the three single deltas is $+0.155$, the bundle sits at
$+0.126$; the residual is $-0.029$ against a spread of $0.61$ ($p=0.84$). The point
estimate is thus slightly sub-additive ($0.82\times$), but statistically indistinguishable
from pure additivity --- the noise dominates. There is no synergy in any case.

\subsection{A Generative Market Model Instead of Rolling Windows}

The second idea comes from the paper selection (Kuo et al., LOB-GAN, see
Chapter~\ref{sec:related-work}): does the agent generalise better if it sees \emph{several
plausible market paths} instead of a single historical one? The obvious route would be
rolling training windows --- and that is \textbf{not buildable} in this protocol. The
environment forces more than $2 \cdot 252 = 504$ observations per environment
(\texttt{assert len(df) > 2*trading\_days}), and \texttt{trading\_days} simultaneously sets
the episode length of 252 steps; the minimum therefore cannot be lowered without breaking
the protocol. The training window, however, has only \textbf{634} usable rows, and the price
series starts exactly at the window's beginning. Two genuine sub-windows would each need
around 550 rows --- impossible. The code now checks this itself and aborts with a clear
message.

Instead, a \textbf{block bootstrap} generates synthetic price series. What matters is
\emph{what} is resampled: price levels would be wrong, because every block seam would create
an artificial jump (a block ends at 100, the next starts at 80 --- minus 20\,\% overnight).
The \emph{daily shape} is therefore bootstrapped:
$\log(\text{Close}_t/\text{Close}_{t-1})$, $\log(\text{Open}_t/\text{Close}_{t-1})$,
$\log(\text{High}_t/\text{Close}_t)$, $\log(\text{Low}_t/\text{Close}_t)$ and
$\log(\text{Volume}_t)$. Blocks of 20 consecutive days are drawn with replacement and the
path is rebuilt from the first real close. Within the blocks the autocorrelation is
preserved, the curve is continuous by construction, and the agent never sees the \emph{real}
training path. Cross-check: the spread of the daily returns is $0.0229$ and $0.0247$ in the
synthetic series against $0.0239$ in the real series.

\subsection{Isolation: the Bootstrap Carries, the Two Passengers Do Not}

In the ``bundle $+$ bootstrap'' row the bootstrap still sits together with the two inert
factors. The isolating run \texttt{nstep} $10$ $+$ bootstrap without \texttt{ent\_coef} and
\texttt{obs\_noise} closes that gap (last row):

\begin{table}[H]
  \centering
  \caption{Paired tests over the 20 seed differences (df$=19$), test margin.}
  \label{tab:update-a3c-paired}
  \small
  \begin{tabular}{lrrr}
    \toprule
    Comparison & $\Delta$ test margin & $t$ & $p$\\
    \midrule
    \texttt{nstep} 10 $-$ baseline            & $+0.107$ & $1.69$ & $0.107$\\
    \texttt{ent\_coef} 0.02 $-$ baseline      & $+0.028$ & $0.47$ & $0.643$\\
    \texttt{obs\_noise} 0.01 $-$ baseline     & $+0.020$ & $0.36$ & $0.720$\\
    bundle $-$ baseline                       & $+0.126$ & $1.98$ & $0.062$\\
    bundle $-$ \texttt{nstep} 10              & $+0.020$ & $0.24$ & $0.810$\\
    \texttt{nstep} 10 $+$ bootstrap $-$ baseline & $+0.200$ & $2.35$ & $\mathbf{0.030}$\\
    \texttt{nstep} 10 $+$ bootstrap $-$ \texttt{nstep} 10 & $+0.093$ & $1.16$ & $0.260$\\
    \texttt{nstep} 10 $+$ bootstrap $-$ bundle $+$ bootstrap & $+0.0002$ & $0.003$ & $0.998$\\
    \bottomrule
  \end{tabular}
\end{table}

\paragraph{The bootstrap effect survives in isolation --- and was never carried by the
passengers.} The isolated increment is $+0.093$ and thus \emph{exactly} the same value as
the contaminated estimate ($+0.093$); the ratio is $1.00$. The effect has therefore neither
shrunk nor disappeared. One caveat: on its own it is \textbf{not significant} ($p=0.26$,
12 of 20 seeds positive, median $+0.138$, trimmed mean $+0.112$) --- the point estimate is
stable, the sample is too small.

\paragraph{\texttt{nstep} 10 $+$ bootstrap is indistinguishable from ``bundle $+$
bootstrap''} ($+0.0002$, $p=0.998$). \texttt{ent\_coef} and \texttt{obs\_noise} are thus
\textbf{finally proven to be passengers}: the two runs differ only in those two parameters
and produce practically the same distribution. (Incidentally a consistency check that the
bootstrap replicates are deterministic.)

\paragraph{A significant edge for the first time.} \texttt{nstep} $10$ $+$ bootstrap is
significantly better than the baseline ($+0.200$, $p=0.030$) and its own test margin is
clearly different from zero ($p=0.0010$, two-sided $0.0019$). This is the \emph{first} arm
in this project to break the 5\,\% threshold in the paired test against the baseline ---
neither \texttt{nstep} alone ($p=0.107$) nor the bundle ($p=0.062$) manages that.

\begin{figure}[H]
  \centering
  \includegraphics[width=\textwidth]{update_2026_a3c.png}
  \caption{One-factor decomposition and block bootstrap, 20 seeds. Left the test margin
  (error bars: 95\,\% CI), right the margin in the training window. The two bootstrap arms
  (dark) have practically zero in-sample edge and the largest out-of-sample one.}
  \label{fig:update-a3c}
\end{figure}

\subsection{What the Overfitting Picture Shows}

Figure~\ref{fig:update-a3c} makes the core visible. The four arms that train on the real
window carry an in-sample margin of between $+6.7$ and $+8.8$ NAV points --- and almost
none of it survives out-of-sample. The two bootstrap arms have an in-sample margin of
$+0.02$, i.e.\ \emph{zero} edge, and at the same time deliver the best test values. They
never saw the historical path; that they do not win in the training window is therefore not
a weakness but the mechanism. An in-sample edge of $+8$ is thus exposed as what it is:
memorisation.

\subsection{Limitations}

\paragraph{20 seeds are a small sample.} The confidence intervals in
Table~\ref{tab:update-a3c-arms} are wide, and the central bootstrap increment is
underpowered at $p=0.26$. What is robust: that the baseline has no edge, that
\texttt{nstep} $10$ explains the bundle, that the two remaining factors do nothing, and
that \texttt{nstep} $10$ $+$ bootstrap beats the baseline significantly. What is not robust
is the claim that the bootstrap addition is significant \emph{by itself}.

\paragraph{Only five different synthetic paths.} With 20 agents and five bootstrap
replicates, groups of four seeds share the same synthetic series. The seeds then vary only
the exploration, not the environment. A run with one path per agent would be the next clean
step.

\paragraph{The in-sample margin of the bootstrap arms is not comparable.} It is measured on
the real window, on which these agents never trained; the row is therefore readable only as
an overfitting indicator, not as a performance measure.

\paragraph{Data and environment.} The limitations from Section~\ref{sec:update-deviations}
apply: observation vector 334 instead of 205, local CSVs instead of \texttt{yfinance}. A3C
uses the same environment as A2C, so the deviations affect both arms equally.

\subsection{Conclusion of the A3C Part}

The baseline confirms the A2C finding on a second algorithm: no demonstrable edge over
buy-and-hold. The decomposition shows that the effect of the obvious ``regularisation''
actually comes from the \textbf{credit assignment} (\texttt{nstep} $5\to10$); entropy bonus
and observation noise are ineffective at this magnitude. The \textbf{block bootstrap}
delivers, as a second lever, the best test value and the first significant edge over the
baseline --- but only in combination. Conceptually this is the more interesting finding: the
agent does not need the real historical path to be level with or better than buy-and-hold on
the test window; many plausible paths are enough.
